\subsection{Oper 奇点与 Bethe 方程}\label{opbethe}


现在我们把具有无单值性奇点的 oper 究竟是怎样一种对象讲清楚（详细得多的讨论见文献~\cite{Frenkel}）。我们将就 $C$ 单连通的情形作这一分析，于是平坦丛 $E\to C$ 的选取不带有模。因此 $C$ 将是 $\R^2$ 或 $\Bbb{CP}^1$，也就是说，它是复 $z$ 平面，添上或不添上无穷远点。假定 $C$ 单连通有两个理由：这是理解 Jones 多项式最相关的情形；并且，把 $E$ 的选取从讨论中排除，会使我们更容易把注意力集中在 oper 及其奇点上。

由于 $C$ 单连通，$C$ 上的平坦丛是平凡的。于是我们可以转到 $\A_z = \A_{\bar z}=0$ 的规范。$E$ 的线子丛 $\L$ 在 $C=\R^2$ 时必然平凡，在 $C=\Bbb{CP}^1$ 时则在略去点 $z=\infty$ 之后平凡。因此它由一个整体截面 $s$ 生成，但我们无法在同时令 $\A_z=\A_{\bar z}=0$ 的情况下把 $s$ 写成简单的形式。于是我们取\begin{equation}\label{nuptriz}s=\begin{pmatrix} P(z)\cr Q(z)\end{pmatrix},\end{equation}
其中 $P$ 与 $Q$ 为多项式。$P$ 与 $Q$ 在相差一个线性变换的意义下唯一确定
\begin{equation}\label{helf}\begin{pmatrix}P\cr Q
\end{pmatrix}\to M\begin{pmatrix}P\cr Q
\end{pmatrix},~~M\in GL(2,\C).\end{equation}

现在 $s\wedge \D_z s$ 化为 $P\partial_zQ-Q\partial_zP$。因此，若多项式 $K(z)=\prod_{a=1}^d(z=z_a)^{k_a}$ 编码了 oper 奇点的位置与荷，则我们希望求解的方程就是
\begin{equation}\label{omicro} P\partial_z Q-Q\partial_zP = K(z),\end{equation}
% [Tutor Note: 原文笔误] 上一段正文中 $K(z)=\prod_{a=1}^d(z=z_a)^{k_a}$ 的因子按上下文应为 $(z-z_a)^{k_a}$；此处依原文逐字保留，未作改动，推导时请以正确的因子为准。
模去 $SL(2,\C)$ 的作用。（取 $K$ 的首项系数为 1 已把 $GL(2,\C)$ 约化为 $SL(2,\C)$。）方便的做法是，通过要求多项式 $Q$ 的次数小于\footnote{后文在存在对称性破缺时，我们将不得不放宽这一条件。} $P$ 的次数（若两个多项式次数相同，则它们的某个线性组合次数更小，我们就把这个组合称为 $Q$），并要求 $Q$ 的首项系数为 1，来固定 $SL(2,\C)$ 的三个生成元中两个的作用，
\begin{equation}\label{zeroes}Q(z)=\prod_{i=1}^q (z-w_i),\end{equation}
其中 $w_i$ 为某些数。这些条件只留下给 $P$ 加上 $Q$ 的倍数的自由度。

我们可以把 (\ref{omicro}) 改写为
\begin{equation}\label{operbo}
\partial_z \frac{P}{Q}  =- \frac{K(z)}{Q^2}
\end{equation}
该式左端在 $Q(z)$ 的零点 $w_i$ 处留数为零。右端也必须留数为零。这给出约束
\begin{equation}\label{bethe}
\sum_a \frac{k_a}{w_i - z_a} = \sum_{j\not=i} \frac{2}{w_i - w_j},~~i=1,\dots,q.\end{equation}
% [Tutor Note: 省略计算] 由 $K/Q^2$ 在单零点 $w_i$ 处（二阶极点）留数为零，即 $\frac{\d}{\d z}(K/q^2)\big|_{w_i}=0$（其中 $Q=(z-w_i)q$，$q(w_i)=Q'(w_i)$），并利用 $K'(w_i)/K(w_i)=\sum_a k_a/(w_i-z_a)$ 以及 $Q''(w_i)/Q'(w_i)=\sum_{j\not=i}2/(w_i-w_j)$，即得 (\ref{bethe})；右端的系数 $2$ 正来自 $Q$ 在 $w_i$ 处的单零点，不要写成 $1$。
反之，给定这些方程的一组解，$K/Q^2$ 的留数为零；于是 $\int K/Q^2\, \d z$ 是有理函数 $P/Q$，而 $P$ 被确定到可加上 $Q$ 的常数倍，这正是预期的未定程度。

若 $C=\R^2$，具有所需无单值性奇点的 oper 就简单地对应于方程 (\ref{bethe}) 的解。对于 $C=\CP^1$，我们还须保证该 oper 在无穷远处没有额外的奇点。结果表明，其条件是多项式 $Q$ 的次数 $q$ 恰为 $K$ 的次数的一半：
\begin{equation}\label{truy} q=\frac{k}{2},~~k=\sum_a k_a. \end{equation}
% [Tutor Note: 易错点] (\ref{truy}) 要求 $k=\sum_a k_a$ 为偶数，这正是后文 (\ref{duflo}) 处 $G^\vee$-不变（$\cmmib w=0$）的条件。对 $C=\R^2$ 并无此要求，故 $C=\R^2$ 与 $C=\CP^1$ 两种情形的 oper 计数不同，比较时切勿混淆。
为判断该 oper 在无穷远处是否有奇点，设 $p$ 为 $P$ 的次数，并定义 $\begin{pmatrix}\tilde P\cr \tilde Q\end{pmatrix}=
z^{-p}\begin{pmatrix}P\cr Q\end{pmatrix}$。我们把 $\tilde P,\,\tilde Q$ 看作 $v=1/z$ 的多项式。$P$ 的次数大于 $Q$ 的次数 $q$ 这一条件蕴含 (\ref{operbo}) 左端 $z$ 的最高次幂不会相消，从而 $p+q=k+1$。由于 $q<p$，一般地有
\begin{equation}\label{ruy} q\leq \frac{k}{2}\end{equation}
该 oper 在 $z=\infty$ 即 $v=0$ 处没有奇点的条件是 $\left(\tilde P\partial_v\tilde Q-\tilde Q\partial_v\tilde P\right)_{v=0}\not=0$，并且不难看出这与 (\ref{truy}) 一致。

为进一步深入，我们需要可积系统理论。相当“神奇”的是，方程 (\ref{bethe}) 是某个可积模型的 Bethe 方程，该模型就是 Gaudin 模型，或 XXX 自旋链的某个大杂质极限。（具有无单值性奇点的 oper 与 Gaudin 模型之间的联系在文献~\cite{Frenkel} 中有综述，其依据是文献~\cite{BF,FFR} 等较早的工作。关于这一点，另见第~\ref{applic}~节。）对 $a=1,\dots,d$，令 $R_a$ 为 $SU(2)$ 的自旋 $j_a=k_a/2$ 表示的一个副本，并令 $\H=\otimes_{a=1}^d R_a$。Gaudin 模型的哈密顿量是 $\H$ 上的一族相互对易的算符
\begin{equation}\label{gaudham}H_a=\sum_{b\not=a}\frac{\vec T_a\cdot \vec T_b}{z_a-z_b}.\end{equation}
这里 $\vec T_a$ 是作用在 $R_a$ 上的 $\frak{su}(2)$ 生成元，而 $\vec T_a\cdot\vec T_b$ 是 $\vec T_a$ 与 $\vec T_b$ 的内积（所用的二次型使得 $\vec T_a\cdot \vec T_a=j_a(j_a+1)$）。实际上，我们在 (\ref{gaudham}) 中写出的是 $G^\vee=SU(2)$ 的 Gaudin 模型的哈密顿量。（我们把这个群记作 $G^\vee$，因为它是本文余下分析中出现的规范群 $G$ 的自然对偶。）对任意 $G^\vee$ 都有相应的 Gaudin 模型，它与 oper 的关系正与我们即将对 $SU(2)$ 描述的关系相同；但若 $G^\vee$ 的秩大于 1，则 $H_a$ 只是一组完备对易哈密顿量的一部分。

由于 Gaudin 哈密顿量相互对易，它们可以同时对角化。此外，由于它们与 $G^\vee$ 的作用对易，其共同本征向量可以按 $G^\vee$ 的不可约表示组织起来。因为有 $G^\vee$ 的作用，要理解 Gaudin 哈密顿量的全部共同本征向量，只需理解那些同时是 $G^\vee$ 作用的最高权向量的共同本征向量。

在 $G^\vee=SU(2)$ 的 Gaudin 模型的 Bethe 拟设理论中，已经证明 Bethe 方程 (\ref{bethe}) 的解对应于那些同时是 $G^\vee$ 作用最高权向量的共同本征向量。在这一对应中，权 $\cmmib w$ 与 $Q$ 的次数 $q$ 由下式相联系：
\begin{equation}\label{duflo} \cmmib w=\frac{k}{2}-q.\end{equation}
特别地，若我们要求 Gaudin 哈密顿量的 $G^\vee$ 不变共同本征向量，则需要 $\cmmib w=0$ 且 $q=k/2$；但正如我们在 (\ref{truy}) 中所见，这正是相应的 oper 能延拓到 $\CP^1$ 上而在无穷远处无奇点的条件。

于是，$\CP^1$ 上具有无单值性奇点的 oper 的数目，与 Gaudin 哈密顿量的 $SU(2)$ 不变共同本征向量的数目相同。但相互对易的 Gaudin 哈密顿量的共同本征向量构成 $\H$ 的一组基，类似地，$G^\vee$ 不变的共同本征向量构成 $\H^{G^\vee}$（即 $\H$ 的 $G^\vee$ 不变部分）的一组基。因此，满足我们所加条件的 oper 的数目恰好是 $\H^{G^\vee}$ 的维数。

这一结果是我们在比较 $G$ 规范场论中 BPS 解的计数与规范群为 $G^\vee$ 的 Chern-Simons 理论时所获得的第一个具体成功。考虑 $\Bbb{CP}^1$ 上的 Chern-Simons 理论，其荷处于表示 $R_a$ 中，$a=1,\dots,d$。我们把这些荷置于点 $z_a\in \Bbb{CP}^1$。在经典极限下，物理态空间就是 $\H=\otimes_a R_a$ 的 $G^\vee$ 不变部分；限制到 $G^\vee$ 不变态就是 Gauss 定律约束。当 Chern-Simons 耦合参数 $\k^\vee$ 取一般值时，这也给出 Chern-Simons 理论物理 Hilbert 空间维数的正确答案。另一方面，若改用对偶描述，以 $G$ 规范场论中与时间无关的解为出发点来构造 Hilbert 空间，则按我们已经给出的论证，这些态应当对应于\footnote{我们在第~\ref{fourdintro}~节解释为什么 $G$ 规范场论中依赖时间的瞬子修正不影响这里的态计数。} $\Bbb{CP}^1$ 上在 $z_a$ 处带荷 $k_a$ 的奇点的 oper。由于这些 oper 的数目与 $\H^{G^\vee}$ 的维数相同，我们至少成功地使两种描述中物理态空间的维数彼此相容。这提示：借助我们在本节中所利用的、向 $\zeta\not=0$ 的形变，第~\ref{methods}~节中基于沿一个方向拉伸纽结的方案确实可以奏效。

若取 $C=\R^2$ 而非 $\CP^1$，态的计数就不那么透明了。从定性上说，Chern-Simons 理论中 $\R^2$ 上的物理态数目多于 $\CP^1$ 上的，因为通量可以逃逸到无穷远，物理态不必完全规范不变。然而，要理解在无穷远处应当施加什么条件相当微妙，而且在 $G^\vee$ 的 Chern-Simons 理论中也难以理解似乎由 oper 得出的结果：$\R^2$ 上的物理态对应于 $\H$ 中的最高权向量。在把对称性破缺纳入考虑（第~\ref{moresymbr}~节）之后，两种描述之间的比较将会更简单。


\subsubsection{与共形场论的关系}\label{relconf}



在得到 Gaudin 模型的过程中，我们所完成的远不止是得到一个与 Jones 多项式的已知构造相呼应的经典解数目。辫群的 Jones 表示可以（见文献~\cite{TK}）描述为 Knizhnik-Zamolodchikov 方程（文献~\cite{KZ}）的单值性。这些方程如下。把进入 Jones 多项式的通常参数 $q$ 写作 $q=\exp(2\pi i/(\k^\vee+2))$。令 $\eusm B$ 为互不相同的 $d$ 元组 $z_1,\dots,z_d\in\C$ 构成的空间，并令 $\H^*$ 为 $\eusm B$ 上以 $\H=\otimes_{a=1}^d R_a$ 为纤维的平凡丛。Knizhnik-Zamolodchikov 方程是 $\H^*$ 的截面 $\Theta$ 所满足的如下方程组：
\begin{equation}\label{kz}\left(\frac{\partial}{\partial z_a}+\frac{H_a}{\k^\vee+2}\right)\Theta=0.
\end{equation}
其中 $H_a$ 即 Gaudin 哈密顿量 (\ref{gaudham})。Knizhnik-Zamolodchikov 方程描述 $\H^*$ 的截面 $\Theta$ 关于某个平坦联络的平行移动，该联络隐含地由 (\ref{kz}) 定义。为验证该联络的平坦性，需利用 $H_a$ 相互对易这一事实以及关系 $\partial H_b/\partial z_a=\partial H_a/\partial z_b$。Knizhnik-Zamolodchikov 方程的解 $\Theta$ 是对称群为 $G^\vee$ 的二维流代数的共形块；它们在二维共形场论中很重要。

由于具有无单值性奇点的 oper 对应于 $\H$ 的一组基，在本文处理 Jones 多项式的方案中，我们最终将用规范场论在丛 $\H^*$ 上构造一个平坦联络；而且，既然我们希望重现辫群的 Jones 表示，这个平坦联络就应当与 Knizhnik-Zamolodchikov 方程所定义的那个规范等价。实际上，这条听起来诱人的路线并不是我们要走的路线。辫群的同一批 Jones 表示在共形场论中还有另一种实现，即用退化场乘积的关联函数的 Virasoro 共形块来实现，而后者将被证明更为有用。




\subsection{再次讨论对称性破缺}\label{moresymbr}

正如第~\ref{symbr}~节那样，我们离开库仑分支原点可以进一步看清一些问题。一如既往，我们的做法是在 $y=\infty$ 附近为 $\vec \phi$ 开启常数且相互对易的期望值。在当前语境下，这意味着联络形式 $\A$ 在无穷远处不为零，而是一个系数为常数的 1-形式；并且这些系数彼此对易。

\subsubsection{“实”对称性破缺}\label{realsym}

首先考虑只有 $\A_y$ 在无穷远处有期望值的情形。这个期望值来自 $\phi_1$ 在无穷远处的值，因此正如 (\ref{gauga}) 那样有
\begin{equation}\label{gaugan}
\A_y = \begin{pmatrix} \ca_1 & 0 \cr 0 & -\ca_1\end{pmatrix},~~y\to\infty,
\end{equation}
其中可取 $\ca_1>0$。

$\A_z$ 与 $\A_{\bar z}$ 在无穷远处不应有指数增长，这一条件告诉我们：在使 $\A_y$ 于 $y\to\infty$ 时形如 (\ref{gaugan}) 的规范中，它们必须是上三角形的。于是，在 $y=\infty$ 附近，实对称性破缺为我们提供了一种把整个平坦联络写成三角形形式的自然方式。

带有结构群到上三角矩阵群（即到 Borel 子群）的协变常值约化的 oper 称为 Miura oper。这一概念在文献~\cite{Frenkel} 中有详细描述。任何无单值性的 oper 丛都可以赋予 Miura oper 结构；事实上，这样做的方式构成一个单参数族。具体地说，若秩 2 平坦丛 $E\to C$ 具有平凡单值性，则与之相伴的 $\Bbb{CP}^1$ 丛（其纤维由 $E$ 的纤维投影化得到）也带有一个无单值性的平坦联络。我们把这个丛记作 $\B$。在某个给定点 $p\in C$ 处任意选取 $\mathcal B$ 的一个截面并作平行移动，就得到 $\mathcal B$ 的一个协变常值截面，它把底层的 oper 变成一个 Miura oper。这一手续引入一个复模——即在起点 $p$ 处 $\B$ 的纤维中选取一点。这意味着无单值性的 Miura oper 依赖于一个复模。（当——如在第~\ref{compsym}~节中那样——我们沿复方向引入对称性破缺时，这个模将会消失，因为届时在无穷远处没有对 $\A_y$ 相对于 $\A_z$ 作规范转动的自由。）

在 Miura oper 的情形下，Bethe 根有一个特别好的解释。为最简单地说明这一点，让我们回到 $C=\R^2$ 的情形并使用 $\A_z=\A_{\bar z}=0$ 的规范。$y\to\infty$ 时的行为挑出秩 2 丛 $E$ 的一个在平行移动下不变的子丛 $\tilde\L$；我们可以称它为平坦子丛。在具有渐近行为 (\ref{gaugan}) 的规范中，$\tilde \L$ 由一个在 $y\to\infty$ 时消失的协变常值截面 $\tilde s$ 生成。经过复规范变换使 $\A_z=\A_{\bar z}=0$ 之后，$\tilde s $ 就简单地成为常数；我们不妨取
\begin{equation}\label{rofor}\tilde s=\begin{pmatrix} 1 \cr 0 \end{pmatrix}.\end{equation}
另一方面，$y\to 0$ 时的行为决定一个全纯（非平坦）子丛 $\L\subset E$，它如前面那样（在 $\A_z=\A_{\bar z}=0$ 的规范中）由
\begin{equation}\label{wagan}s=\begin{pmatrix} P \cr Q\end{pmatrix}.\end{equation}
生成。Miura 结构 $\tilde s$ 的选取提供了一种挑出 $s$ 各分量的自然线性组合的方式，即
\begin{equation}\label{agan}\tilde s \wedge s=Q. \end{equation}
这样确定的 $Q$ 未必是我们在第~\ref{opbethe}~节中所用的那个：在那里我们取 $Q$ 为 $s$ 各分量中次数较小的线性组合；而现在 $Q$ 简单地由 Miura 结构决定。按我们的新选择，$Q$ 的零点有一个简单的解释：它们是 $s$ 成为 $\tilde s$ 的倍数的那些点。换言之，由 $y\to\infty$ 的行为所决定的子丛，一般不同于由 $y\to 0$ 的行为所决定的子丛。$Q$ 的零点——称为 Bethe 根——恰恰是二者重合的那些点。换一种说法：Bethe 根是这样的 $z$ 值，在它们处存在一个 $\D_ys= 0$ 的解，且该解在 $y\to 0$ 与 $y\to\infty$ 时都消失。正如我们在第~\ref{symbr}~节中所讨论的，在 Bogomolny 方程的语境下，人们会说在那些 $z$ 值（以及某些 $y$ 值）处存在光滑的 BPS 磁单极子。
% [Tutor Note: 物理图像] 来自 $y\to\infty$ 的平坦子丛 $\tilde\L$ 与来自 $y\to0$ 的全纯子丛 $\L$ 一般不同；Bethe 根就是二者重合的 $z$ 值，那里存在一个沿 $y$ 方向、在两端都消失的解，即一个光滑的 BPS 磁单极子。Bethe 方程的代数条件由此获得清晰的几何意义。


作为一个具体例子，假设根本不存在 't Hooft 算符。平坦丛 $E\to \R^2$ 完全平凡，且在同构意义下具有唯一的 oper 结构：
\begin{equation}\label{only} s=\begin{pmatrix} z \cr 1\end{pmatrix}.\end{equation}
在没有对称性破缺时，我们关于 $Q$ 是 $s$ 各分量中次数较小的线性组合的约定导致 $Q=1$，从而（在给 $P$ 加上 $Q$ 的倍数并重新标度的自由度之内）有 $P=z$。

在存在实对称性破缺时，Miura 结构给出 $Q$ 的一个特殊选择 (\ref{agan})，它没有理由为常数。若 $Q$ 不是常数，则通过给 $P$ 加上 $Q$ 的倍数并重新标度，可以取 $P=1$，于是
\begin{equation}\label{onl} s=\begin{pmatrix}P \cr Q \end{pmatrix}=\begin{pmatrix}1 \cr z-w
\end{pmatrix},\end{equation}
其中 $w$ 为某个数。多项式 $K=PQ'-QP'$ 为 1，这与不存在任何 't Hooft 算符相一致。我们发现，在存在实对称性破缺时，全纯数据对应于一个依赖于 $w\in \R^2$ 的选取的单参数族解。这是一个不含 't Hooft 算符、只有一个 Bethe 根的解。我们称它为裸 Miura oper。在适当的语境中，当把它提升回四维时，我们稍后会将这个解解释为一个定域在 $z=w$ 处以及某个依赖于 $\zeta$ 的 $y$ 值处的“弦”。
% [Tutor Note: 约定差异] 这里的 “裸 Miura oper” 含有一个复模 $w$，它与无单值性 oper 的一般 Miura 结构所带的模是同一回事；在后文引入复方向对称性破缺（第~\ref{compsym}~节）后该模消失，所以此处不要把它误当作物理自由度来计数。

一般地，$K$ 的次数至多比 $Q$ 的次数小 1（此界恰在 $P=1$ 时达到），因此 Bethe 根的数目至多比 $K$ 的次数 $k=\sum_ak_a$ 多 1。

